\newif\ifcontest\contestfalse
\documentclass[a4paper,zihao=-4,fontset=fandol]{ctexart}
\usepackage[left=22.5mm,right=22.5mm,top=25mm,bottom=25mm,footskip=12mm]{geometry}
\usepackage{amsmath,amssymb,bm,graphicx,booktabs,tabularx,longtable,array,float}
\usepackage{caption,fancyvrb,algorithm,algorithmic,hyperref}
\hypersetup{hidelinks,pdftitle={低空湍流监测与航路优化：可编辑写作模板},pdfauthor={},pdfsubject={2025 D题历史题写作演练}}
\setmainfont{texgyretermes}[Extension=.otf,UprightFont=*-regular,BoldFont=*-bold,ItalicFont=*-italic,BoldItalicFont=*-bolditalic]
\IfFileExists{build/fonts.tex}{\input{build/fonts}}{}
\ifcontest
  \linespread{1.0}
\else
  \geometry{top=30mm}
  \linespread{1.38}
\fi
\setlength{\parindent}{2em}
\setlength{\parskip}{0pt}
\setlength{\emergencystretch}{2em}
\ctexset{section={format=\centering\heiti\bfseries\zihao{4},beforeskip=18bp,afterskip=12.4bp},
subsection={format=\bfseries\zihao{-4},beforeskip=12bp,afterskip=6.8bp},
subsubsection={format=\bfseries\zihao{-4},beforeskip=10bp,afterskip=6bp}}
\pagestyle{plain}
\renewcommand{\thefigure}{\thesection.\arabic{figure}}
\renewcommand{\thetable}{\thesection.\arabic{table}}
\renewcommand{\theequation}{\thesection.\arabic{equation}}
\captionsetup{font=normalsize,labelsep=quad,skip=8bp}
\floatname{algorithm}{算法}
\renewcommand{\algorithmicrequire}{\textbf{输入：}}
\renewcommand{\algorithmicensure}{\textbf{输出：}}
\newcolumntype{Y}{>{\raggedright\arraybackslash}X}
\newcommand{\pending}{\textbf{待计算}}
\newcommand{\slot}[1]{\textbf{【#1】}}
\newcommand{\E}{\mathbb{E}}
\newcommand{\R}{\mathbb{R}}
\newcommand{\samplefigure}[3]{%
  \begin{figure}[H]\centering
  \includegraphics[width=0.91\linewidth,height=0.43\textheight,keepaspectratio]{figures/#1.pdf}
  \caption{#2（示意占位图）}\label{#3}\end{figure}}
\newcounter{extrafigure}
\renewcommand{\theextrafigure}{S.\arabic{extrafigure}}
\newcommand{\auxfigure}[3]{%
  \begin{figure}[H]\centering
  \includegraphics[width=0.91\linewidth,height=0.40\textheight,keepaspectratio]{figures/#1.pdf}
  \refstepcounter{extrafigure}\caption*{图\theextrafigure\quad #2（补充示意占位图）}\label{#3}\end{figure}}
\newcounter{extratable}
\renewcommand{\theextratable}{S.\arabic{extratable}}
\newcommand{\auxtablecaption}[1]{\refstepcounter{extratable}\caption*{表\theextratable\quad #1}}
\newcommand{\aitool}{\slot{工具名称}}
\newcommand{\aimodel}{\slot{版本或型号}}
\newcommand{\aideveloper}{\slot{开发机构}}
\newcommand{\aireleased}{\slot{版本发布日期}}
\newcommand{\AIDataNote}{\ifcontest\par\noindent
  AI 辅助数据分析标注：本处涉及的工具为\aitool，版本或型号为\aimodel，
  开发机构为\aideveloper，版本发布日期为\aireleased；采用内容为\slot{实际采用内容}，
  队伍复核方式为\slot{具体复核方法与结果}。\par\fi}
\newcommand{\modelcard}[4]{\subsection{#1}#2\par\textbf{基本形式：}#3\par\textbf{选择与检验：}#4\par}

\hypersetup{pdftitle={常见模型备选与替换段落}}
\begin{document}
\begin{center}{\heiti\zihao{3}常见模型备选与替换段落}\end{center}
本手册与论文正文分开使用。每个条目给出可直接改写的建模段落、基本形式、适用输入和检验重点。最终论文应保留实际采用的方法，补充真实参数、来源和结果；模型名称的数量不作为方案质量的依据。
\section{回归、预测与特征提取}
\modelcard{线性回归与岭回归}{针对解释变量与目标量之间的近似线性关系，首先建立线性基线。考虑到多个观测特征之间可能存在相关性，在残差平方和中加入二范数惩罚，以控制系数波动。该方法用于双模型标定、分层偏差订正及多变量预测。}{\[\min_{\beta_0,\bm\beta}\frac1n\|\bm y-\beta_0\bm1-X\bm\beta\|_2^2+\lambda\|\bm\beta\|_2^2.\]}{输入为同一尺度的特征矩阵和连续目标，标准化只在训练数据上拟合。比较残差结构、时间块误差与系数稳定性；截距不惩罚。}
\modelcard{LASSO与弹性网络}{当候选变量较多时，通过稀疏惩罚筛选与目标相关的变量。弹性网络在稀疏性之外保留二范数项，以缓和高度相关变量之间的替代关系。}{\[\min_{\bm\beta}\frac1{2n}\|\bm y-X\bm\beta\|^2+\lambda_1\|\bm\beta\|_1+\lambda_2\|\bm\beta\|_2^2.\]}{用于可解释变量筛选。检查不同重采样划分下的入选频率，避免把单次入选直接解释为因果机制。}
\modelcard{多项式回归与样条}{当散点关系表现出平滑弯曲趋势时，引入低阶多项式或局部样条基函数，使模型在保持可解释结构的同时描述非线性。}{\[\widehat y=\beta_0+\sum_{j=1}^{K}\beta_jB_j(x).\]}{优先比较二次项和少量节点样条，检查边界外推、峰值振荡和自由度。三小时历史序列中避免用高阶多项式追踪噪声。}
\modelcard{随机森林回归}{针对变量之间的非线性作用及交互关系，采用多棵回归树对不同样本子集进行学习，再对树输出取平均。分裂规则刻画局部关系，叶节点样本数控制拟合的细致程度。}{\[\widehat y(\bm x)=B^{-1}\sum_{b=1}^{B}h_b(\bm x).\]}{适合模型 b 标定与低维非线性外推。调节深度、叶节点样本数和树数，采用时间或站点分组验证；随机行划分可能高估性能。}
\modelcard{梯度提升树}{为逐步修正已有模型的残差，采用加法模型迭代构造回归函数。每轮弱学习器拟合当前损失的负梯度，以学习率控制更新幅度。}{\[F_m(\bm x)=F_{m-1}(\bm x)+\eta h_m(\bm x).\]}{适合表格数据和复杂交互。用独立验证块早停，比较持续性与岭回归基线；训练误差很低时仍需检查测试误差和空间外推。}
\modelcard{支持向量回归}{在样本规模适中且映射较平滑的情形下，采用核函数描述特征空间中的非线性关系，并利用不敏感损失降低小幅误差对拟合的影响。}{\[\min_{f}\tfrac12\|f\|_{\mathcal H}^2+C\sum_i\max\{0,|y_i-f(\bm x_i)|-\epsilon\}.\]}{输入特征需要标准化，核宽度与惩罚参数联合验证。样本较大时记录内存和求解成本。}
\modelcard{高斯过程回归}{将未知函数建模为带有均值和协方差核的随机函数，利用观测之间的相关关系同时估计均值和预测方差。}{\[f\sim\mathcal{GP}(m,k),\qquad \mu_*=m_*+K_{*X}(K_{XX}+\sigma^2I)^{-1}(\bm y-\bm m).\]}{适合中小样本和空间不确定性分析。核假设需要验证，预测区间应检查覆盖率；大样本使用稀疏近似并说明诱导点选择。}
\modelcard{ARIMA与带外生变量的ARIMAX}{当目标具有明显的序列相关性时，先通过差分处理趋势，再以历史目标和历史误差解释当前变化。可获得的外生变量作为额外输入进入回归项。}{\[\phi(B)(1-B)^dy_t=\theta(B)\epsilon_t+\bm\gamma^\mathsf T\bm x_t.\]}{需要足够长且频率一致的历史序列。检验残差自相关、预测区间和滚动误差；未来外生变量不可得时，需明确其预测来源。}
\modelcard{指数平滑与持续性基线}{短时预测先以当前状态持续不变作为基线，再通过指数平滑降低观测噪声。参数控制近期信息与历史信息之间的权重。}{\[\ell_t=\alpha y_t+(1-\alpha)\ell_{t-1},\qquad \widehat y_{t+h}=\ell_t.\]}{计算成本低，适合所有时序候选的基础对照。该模型本身为线性平滑，题面要求非线性模型 e 时应另提供非线性主模型。}
\modelcard{LSTM、TCN与时空图网络}{当具备连续且足量的训练序列时，可用时间网络学习多步依赖，再用空间图模块传递邻近位置的信息。预测窗口、输入窗口和目标时效在训练前固定。}{\[\widehat{\bm y}_{t+1:t+h}=f_\Theta(\bm X_{t-L+1:t},A).\]}{适合存在额外合规历史样本的情况。报告数据覆盖、参数规模、随机种子、早停与时间边界；三小时序列不自动提供足够独立样本。}
\modelcard{主成分分析与POD降维}{针对高维网格变量，利用主要变化方向表示场状态，使后续状态空间建模在低维系数上进行。}{\[X_c=U\Sigma V^\mathsf T,\qquad \bm I_t\approx\overline{\bm I}+V_r\bm a_t.\]}{主成分在训练部分拟合。除累计方差外，还需检查重构峰值误差与下游航路排序，避免低方差局地危险信息被舍弃。}
\section{空间重构与动态估计}
\modelcard{各向异性IDW}{考虑水平与垂直相关尺度不同，采用方向缩放后的距离计算观测权重，并融合质量与时间衰减。}{\[\widehat I(\bm r)=\frac{\sum_iq_i(d_i^2+\epsilon)^{-p/2}I_i}{\sum_iq_i(d_i^2+\epsilon)^{-p/2}}.\]}{用于模型 c 基线。检查无邻居、重复点、零距离和留站误差；距离权重本身不提供概率置信区间。}
\modelcard{普通克里金}{对具有可估计空间相关结构的指标，利用半方差函数构造无偏线性估计，并最小化预测误差方差。}{\[\begin{bmatrix}\Gamma&\bm1\\\bm1^\mathsf T&0\end{bmatrix}\begin{bmatrix}\bm w\\\mu\end{bmatrix}=\begin{bmatrix}\bm\gamma_*\\1\end{bmatrix},\qquad \widehat I=\bm w^\mathsf T\bm I.\]}{需要拟合有效半方差模型，注意各向异性与设备偏差。用空间块或留站验证检查预测误差与克里金方差的校准。}
\modelcard{三维变分融合}{通过观测算子将网格场映射到各设备的采样位置，再将观测拟合、背景约束和平滑要求合并为一个目标。}{\[J=\|H\bm I-\bm y\|_{R^{-1}}^2+\lambda_b\|\bm I-\bm I_b\|^2+\lambda_s\|L\bm I\|^2.\]}{适合多源空间重构。检查算子单位、正定条件、边界规则与正则敏感性；含有界约束时使用匹配的优化算法。}
\modelcard{卡尔曼滤波、扩展卡尔曼与集合卡尔曼}{将状态演化和观测误差分别建模，通过预测与新息更新递推估计状态。线性高斯模型使用卡尔曼滤波，非线性问题采用局部线性化或集合近似。}{\[\bm s_t=F(\bm s_{t-1})+\bm\eta_t,\qquad\bm y_t=H(\bm s_t)+\bm\nu_t.\]}{适合连续观测场的平滑和短时预测。检查创新序列、协方差稳定性与区间覆盖；高维场需要局地化或降维。}
\modelcard{半拉格朗日平流}{在短时间内近似由风场输送风险结构，反向追踪网格点的来源，并在历史场上插值得到未来状态。}{\[I_{t+\Delta t}(\bm r)=\mathcal I[I_t](\bm r-\bm U\Delta t).\]}{适合平移占主导的临近预报基线。检查边界、速度插值、数值耗散及局地生消；标准插值不自动满足守恒。}
\modelcard{反应-平流-扩散模型}{当空间输送、局地扩散和生消均有物理依据时，将场演化写为偏微分方程，再由观测标定系数。}{\[\partial_tI+\bm U\cdot\nabla I=\nabla\cdot(K\nabla I)+S(I,\bm r,t).\]}{需要边界条件、初值和可识别的源项。检验量纲、数值稳定性、非负性与网格收敛，避免用不可辨识参数补偿所有误差。}
\section{优化、决策与验证}
\modelcard{Dijkstra与A*}{将可行飞行状态表示为图节点，将边上的累计风险表示为非负代价，通过最短路算法求解路径。A* 在已知下界时利用启发函数减少搜索。}{\[C(\pi)=\sum_{e\in\pi}c_e,\qquad f(n)=g(n)+h(n),\quad h(n)\le C^*(n,\mathrm{goal}).\]}{主任务只考虑湍流时使用风险代价。无正风险下界时取零启发；动态场使用含时间的状态，检查完整边段的空域可行性。}
\modelcard{线性规划与混合整数规划}{当决策量具有线性约束且包含路径、分配或开关变量时，将问题写为统一的约束优化形式。整数变量用于表示离散选择。}{\[\min_{\bm x}\bm c^\mathsf T\bm x,\quad A\bm x\le\bm b,\quad x_j\in\mathbb Z\ (j\in\mathcal J).\]}{适合资源分配、带组合约束路径及其他华为杯题型。优先采用成熟求解器，报告可行性残差、最优界、gap 与超时状态。}
\modelcard{动态规划}{当问题可以分解为阶段状态与决策，且未来代价由当前状态充分表示时，采用 Bellman 递推求解。}{\[V_t(s)=\min_{a\in A(s)}\{c_t(s,a)+V_{t+1}(f(s,a))\}.\]}{适合有限时间路径和库存决策。状态必须保留影响未来约束的变量，检查终端条件与状态离散误差。}
\modelcard{NSGA-II与多目标Pareto分析}{当多个目标难以事先合并时，生成非支配候选集合，再比较不同偏好下的方案。采用固定预算和多个随机种子评估搜索稳定性。}{\[\min_{\bm x\in\mathcal F}(f_1(\bm x),\ldots,f_m(\bm x)).\]}{适用于明确的多目标扩展，不替代本题仅湍流条件的主目标。检查约束可行性、Pareto 支配关系与简单基线。}
\modelcard{鲁棒优化、机会约束与CVaR}{当输入或预报存在不确定性时，将多个情景下的代价纳入决策，约束高风险事件或控制尾部损失。}{\[\min_\pi\E[C(\pi)]+\lambda\operatorname{CVaR}_\alpha(C),\qquad P(g(\pi,\xi)>0)\le\epsilon.\]}{情景来源应与实际不确定性一致。样本满足率需要独立验证和抽样误差分析，不能直接等同总体概率保证。}
\modelcard{AHP、熵权与TOPSIS}{当题目要求多指标综合评价时，先确定指标方向与尺度，再依据专家判断或样本离散程度给出权重，通过理想点距离形成排序。}{\[C_i=\frac{D_i^-}{D_i^++D_i^-},\qquad D_i^\pm=\sqrt{\sum_jw_j(z_{ij}-z_j^\pm)^2}.\]}{适合评价排序类题目。核对成本型指标、AHP 一致性及权重敏感性；离散程度较高不代表指标更重要。}
\modelcard{聚类、DBSCAN与异常检测}{当缺少类别标签时，依据特征距离识别样本群组或稀疏点，辅助描述不同气象状态及质量异常候选。}{\[\min_{\mu_1,\ldots,\mu_K}\sum_i\min_k\|\bm x_i-\mu_k\|^2.\]}{K-means 适合近似球状簇，DBSCAN 适合密度结构。先做单位与尺度处理，异常点仍需物理检查；聚类结果不直接等同湍流等级。}
\modelcard{Bootstrap与蒙特卡洛敏感性分析}{对时间块、站点或模型参数进行重复采样，重新执行计算链，观察结果指标和路径选择的变化。}{\[\widehat{\mathrm{CI}}=[Q_{\alpha/2}(\widehat\theta^{(b)}),Q_{1-\alpha/2}(\widehat\theta^{(b)})].\]}{时间与空间相关数据使用对应分块采样。记录采样单元、重复次数和随机种子，并分别报告测量、参数与模型结构的不确定性。}
\begin{table}[H]\centering\caption{模型替换时的最小比较集合}
\begin{tabularx}{\linewidth}{lYYY}\toprule
任务 & 基线 & 主要候选 & 比较内容\\\midrule
回归标定 & 岭回归 & 随机森林、核方法 & 分组误差、单位与输入边界\\
空间重构 & IDW & 克里金、变分融合 & 留站误差、覆盖与区间\\
时间预测 & 持续性 & 非线性回归、状态空间 & 分时效误差、时间泄漏\\
路径优化 & Dijkstra & A*、动态规划、MIP & 可行性、同目标值与运行成本\\\bottomrule
\end{tabularx}\end{table}
\samplefigure{modelselection}{数据条件与候选模型的选择关系}{fig:modelselection}
\end{document}
